\section{Введение}

Генерация, дополненная поиском, (Retrieval-Augmented Generation, \gls{RAG}) \cite{gao2023retrieval}
представляет собой подход в задачах обработки естественного языка, сочетающий
информационный поиск и генерацию текста. Вместо того чтобы опираться только на знания,
заложенные в параметрах языковой модели, RAG-система сначала извлекает внешние документы,
а затем использует их при формировании ответа. Такой подход особенно востребован в
прикладных задачах, где важны актуальность ответа, опора на источники и возможность
работать с большими корпусами документов.

Однако повышение качества отдельных компонентов RAG-системы не всегда приводит к улучшению
ее работы в целом. Более релевантное извлечение документов само по себе не гарантирует
формирования корректного ответа, а увеличение полноты генерации может сопровождаться
снижением фактической точности, некорректной атрибуцией источников или появлением
галлюцинаций. В связи с этим оценка RAG-систем требует не только раздельного анализа
компонентов поиска и генерации, но и подхода, позволяющего учитывать их взаимодействие
и влияние на качество итогового ответа.


В последние годы для этой задачи были предложены как LLM-ориентированные фреймворки \gls{Фреймворк}
оценки, так и обучаемые модели-оценщики. Наиболее известными примерами являются
RAGAS \cite{es2024ragas}, RAGBench \cite{friel2024ragbench}, Luna \cite{belyi2024luna},
ARES \cite{saad2023ares}, RAGTruth \cite{niu2023ragtruth} и
AttributionBench \cite{li2024attributionbench}. Эти работы предлагают разные способы
измерять качество ответа, полноту использования контекста и степень его согласованности
с источниками. Среди них особое место занимает RAGBench, поскольку он объединяет крупный
открытый корпус и интерпретируемую схему разметки \gls{TRACe}. По результатам
оригинальной статьи~\cite{friel2024ragbench}, компактный дообученный классификатор
может превосходить большую языковую модель без дообучения~\cite{brown2020language}
по стабильности и точности на англоязычных поддатасетах. Этот результат относится
к оригинальной работе и в настоящей диссертации не воспроизводится в виде прямого
сравнения по единой метрике на одном корпусе, а используется как методологический
ориентир.

Для русского языка ситуация остается существенно менее изученной. Несмотря на рост числа
русскоязычных RAG-приложений, отсутствует общепринятый открытый бенчмарк \gls{Бенчмарк}, совместимый с
TRACe-подходом, а также систематический анализ того, как русскоязычные и мультиязычные
базовые модели ведут себя в задаче токен-уровневой оценки RAG. Это затрудняет как
сравнение систем между собой, так и перенос результатов англоязычных работ на русский
язык.

Настоящая работа посвящена воспроизводимому переносу 
подхода RAGBench/TRACe на русский язык. 
В работе создается русскоязычный бенчмарк, 
обучается семейство моделей-оценщиков и проводится 
серия экспериментов, направленных на сравнение архитектур, 
а также на анализ зависимости метрик от 
длины контекста, схемы весов функции потерь и языка.

\subsection*{Цель и задачи работы}

Цель данной диссертации состоит в разработке и экспериментальной проверке воспроизводимого
подхода к оценке русскоязычных RAG-систем, включающего русскоязычный бенчмарк, семейство
моделей-оценщиков и средства их практического применения.

Для достижения этой цели решаются следующие задачи:

\begin{enumerate}
  \item построить и опубликовать русскоязычный бенчмарк RAGBench-RU на основе 7
  поддатасетов исходного RAGBench, переведенных с помощью LLM;
  \item воспроизвести схему оценки без дообучения на Qwen2.5-72B по аналогии
  с протоколом исходного RAGBench на тестовых поддатасетах англоязычного
  корпуса (см.~приложение~А.3) и использовать её как опорный ориентир;
  прямое сравнение с классификатором по единой метрике и на одном корпусе
  не проводится;
  \item обучить семейство токен-уровневых моделей-оценщиков TRACe, варьируя базовую
  модель, длину контекста, веса функции потерь и архитектуру классификационной головы;
  \item провести сравнительный анализ полученных конфигураций: исследовать
  зависимость метрик от длины контекста, влияние схемы весов функции потерь и
  архитектуры классификационной головы, сравнить качество моделей, обученных
  на русскоязычной и англоязычной версиях корпуса, а также сопоставить
  обученный классификатор с zero-shot LLM-оценщиком;
  \item реализовать интерактивный сервис, позволяющий применять обученные модели к
  произвольным тройкам <<вопрос -- контекст -- ответ>>.
\end{enumerate}

\subsection*{Положения, выносимые на защиту}

\begin{enumerate}
  \item Построен и опубликован русскоязычный бенчмарк RAGBench-RU
  (HuggingFace \url{https://huggingface.co/datasets/CMCenjoyer/ragbench-ru}),
  полученный систематическим переводом 7 поддатасетов исходного RAGBench с сохранением
  структуры разметки TRACe.
  \item Семейство русскоязычных токен-уровневых моделей-оценщиков позволяет с высокой
  точностью предсказывать метки relevance, utilization и adherence; лучшая конфигурация
  достигает среднего Test~F1 = 0.858 на объединенном тестовом корпусе.
  \item Воспроизведена и количественно охарактеризована схема оценки без
  дообучения на базе большой языковой модели Qwen2.5-72B на исходном
  англоязычном RAGBench: средняя доля совпадений по предложениям составляет
  $0.70$ для релевантности и $0.83$ для использования контекста, а также
  выявлена устойчивая склонность к избыточной разметке (в среднем около
  $1.28$ лишних релевантных ключей и $0.73$ лишних использованных ключей на
  пример), наиболее выраженная в юридическом и научном поддатасетах. Эти
  значения используются в работе как опорный ориентир и согласуются с
  ключевой идеей оригинальной статьи RAGBench о практической применимости
  локального классификатора как более дешёвой замены большой языковой модели;
  при этом прямое сопоставление двух методов по единой метрике на уровне
  токенов и на одном корпусе в настоящей работе не проводится и относится к
  направлениям дальнейшего исследования.
  \item На исследованной выборке конфигураций простая классификационная голова
  систематически оказывается устойчивее сложной, а оптимальная длина контекста для
  лучших моделей лежит в диапазоне порядка 1024--4096 токенов с максимумом при 2048.
  \item Разработан прикладной Streamlit-сервис, позволяющий использовать обученные
  оценщики для одиночной и пакетной оценки новых примеров без повторного обучения.
\end{enumerate}

\subsection*{Научная новизна и практическая значимость}

\textbf{Научная новизна.} Насколько известно автору, данная работа является 
первой
систематической попыткой переноса подхода RAGBench/TRACe на русский язык, 
объединяющих публичный бенчмарк, обученные модели-оценщики и анализ  
факторов, влияющих на качество токен-уровневой оценки. В отличие от подходов,
 основанных преимущественно на
вызовах LLM в стиле RAGAS или на закрытых внутренних метриках отдельных команд, в работе
предлагается воспроизводимая исследовательская схема, доступная для дальнейшего сравнения
и развития.

\textbf{Практическая значимость.} Полученные результаты позволяют:
(а)~сравнивать русскоязычные RAG-системы в единой методологии;
(б)~использовать обученные оценщики как локальную и существенно более дешевую альтернативу
LLM-оценке в офлайн-сценариях;
(в)~применять разработанные модели к новым данным через открытый обучающий конвейер и
интерактивный сервис. При этом вопросы полноценной CI-интеграции и сервисного
развертывания рассматриваются как отдельное направление дальнейшей работы.

\subsection*{Апробация результатов}

Основные результаты настоящей работы были представлены и обсуждены
на научной конференции <<Ломоносовские чтения~- 2026>>
(секция <<Вычислительная математика и кибернетика>>),
МГУ имени М.\,В.\,Ломоносова, 23~марта~- 5~апреля 2026~г.
По теме настоящей работы опубликованы тезисы доклада~\cite{moiseev2026ragru}.

\subsection*{Структура работы}

Раздел~\ref{sec:related_work} содержит обзор существующих подходов к оценке RAG-систем и
обосновывает выбор RAGBench/TRACe в качестве методологической основы работы.
В разделе~\ref{sec:problem} формализуется постановка задачи.
Раздел~\ref{sec:dataset} описывает построение корпуса RAGBench-RU, конвейер перевода и
контроль качества.
В разделе~\ref{sec:model} излагается архитектура модели-оценщика, варианты
классификационной головы и процедура подбора порогов.
Раздел~\ref{sec:experiments} посвящен экспериментальному анализу: сравнению моделей,
длин контекста, весов функции потерь, языков разметки и LLM-оценки.
В разделе~\ref{sec:service} описан интерактивный сервис применения обученных моделей.
В разделе~\ref{sec:conclusion} подводятся итоги, обсуждаются ограничения и направления
дальнейшей работы.
